% !PRINTED 701
\weberpage{701}
%<<CONT
. A differential whose lower polygon consists of only two simple points is called a differential of the third kind. We prove the theorem:

4. \gesperrt{Every differential $dJ$ can be composed linearly and with constant coefficients out of differentials of the first, second and third kind.}

This theorem can also be stated, in accordance with (4), by saying that the class $A$ has a basis
\weq{6}{\mathfrak{A}_1, \mathfrak{A}_2', \ldots, \mathfrak{A}_{b+p-1}}
of such a kind that the quotients
\weq{7}{dJ_r = \frac{\mathfrak{A}_r}{\mathfrak{B}}}
are differentials of one of the three kinds.

In order to prove it by complete induction, we assume that for some $\mathfrak{B}$ such a basis has been found, and we seek one of the same sort for the lower polygon $\mathfrak{P}\mathfrak{B}$, where $\mathfrak{P}$ is an arbitrary point. Such a basis we may assume in the form:
\weq{8}{\mathfrak{P}\mathfrak{A}_1, \mathfrak{P}\mathfrak{A}_2, \ldots, \mathfrak{P}\mathfrak{A}_{b+p-1}, \mathfrak{A}_{b+p},}
in which $\mathfrak{A}_{b+p}$ must be a polygon equivalent to $\mathfrak{W}\mathfrak{B}\mathfrak{P}$ that does not contain the point $\mathfrak{P}$.

If the point $\mathfrak{P}$ divides $\mathfrak{B}$, we put:
\[
\mathfrak{B} = \mathfrak{M}\mathfrak{P}^m,
\]
so that $\mathfrak{P}$ no longer divides $\mathfrak{M}$. The polygon $\mathfrak{W}\mathfrak{P}^{m+1}$ then belongs to a proper class, in which therefore a polygon $\mathfrak{N}$ not divisible by $\mathfrak{P}$ exists.

We then put:
\[
\mathfrak{A}_{b+p} = \mathfrak{M}\mathfrak{N},
\]
and from this we obtain a differential of the second kind:
\weq{9}{dJ_{b+p} = \frac{\mathfrak{A}_{b+p}}{\mathfrak{B}\mathfrak{P}} = \frac{\mathfrak{N}}{\mathfrak{P}^{m+1}}.}

If $\mathfrak{P}$ does not divide $\mathfrak{B}$, then one takes a point $\mathfrak{P}_0$ which divides $\mathfrak{B}$. Then $\mathfrak{W}\mathfrak{P}_0\mathfrak{P}$ belongs to a proper class and contains a polygon $\mathfrak{N}$ not divisible by $\mathfrak{P}$. If then
\[
\mathfrak{B} = \mathfrak{P}_0\mathfrak{M}, \quad \mathfrak{A}_{b+p} = \mathfrak{M}\mathfrak{N},
\]
then
\weq{10}{dJ_{b+p} = \frac{\mathfrak{A}_{b+p}}{\mathfrak{B}\mathfrak{P}} = \frac{\mathfrak{N}}{\mathfrak{P}_0\mathfrak{P}}}
is
%CONT>>
